% Source: https://artofproblemsolving.com/community/c3h2518771_rj5303707_birthday_problems
% Stable identity: aops-topic-2518771-q01
Let the legs be $a$ and $b$. Rotating about a leg gives a cone whose height is that leg and whose base radius is the other leg. Rotating about $a$ gives volume $\dfrac13\pi b^2a$, and rotating about $b$ gives $\dfrac13\pi a^2b$. So
\[b^2a=112,\qquad a^2b=196.\]
Multiplying, $(ab)^3=112\cdot196=21952=28^3$, so $ab=28$. Then $b=\dfrac{112}{28}=4$ and $a=\dfrac{196}{28}=7$, and the hypotenuse is
\[\sqrt{7^2+4^2}=\boxed{\sqrt{65}}.\]

\begin{asy}
usepackage("amsmath");
size(4cm);
defaultpen(fontsize(10pt));
real a=7, r=4;
pair O=(0,0), T=(0,a), Rr=(r,0), L=(-r,0);
fill(T--Rr--O--cycle,rgb(0.85,0.9,1));
draw(T--Rr,linewidth(0.9)); draw(T--L,dashed);
draw(shift(0,0)*xscale(r)*yscale(0.9)*arc((0,0),1,180,360),linewidth(0.8));
draw(shift(0,0)*xscale(r)*yscale(0.9)*arc((0,0),1,0,180),dashed);
draw(O--T,linewidth(0.9)); draw(O--Rr,linewidth(0.9));
draw((0,0.45)--(0.45,0.45)--(0.45,0));
label("$7$",(0,a/2),W); draw((0,-1.3)--(r,-1.3),Arrows(4)); label("$4$",(r/2,-1.3),S); label("$\sqrt{65}$",(T+Rr)/2,NE);
dot(T);
\end{asy}

(The sketch shows the rotation about the leg $a=7$: a cone of radius $4$ and height $7$, of volume $\dfrac13\pi\cdot16\cdot7=\dfrac{112\pi}{3}$.)
